% ================================================================
%  §3  Univariate Optimization
% ================================================================
\section{Univariate Optimization}\label{sec:univariate}

Linear and quadratic functions have closed-form optima.
For anything more complex --- a transcendental objective, a polynomial of high degree, a black-box simulation --- we need iterative algorithms.
This section develops the key univariate methods, from slow and robust to fast and fragile. At each step, we trade a stronger assumption on $f$ for a faster algorithm. The section ends with the question: what if the local optimum we find is not the global one?


% ────────────────────────────────────────────────────────────────
\subsection{Lipschitz Continuity: Making Optimization Possible}\label{ssec:lipschitz}

\Cref{ssec:why-hard} showed that optimization is impossible for arbitrary functions --- isolated minima can hide in arbitrarily narrow spikes.
The cure: impose a ``speed limit'' on how fast $f$ can change.

\begin{defn}[Lipschitz continuous function]\label{def:lipschitz}
	A function $f$ is \textit{Lipschitz continuous} with constant $L > 0$ on a set $X$ if
	\begin{equation}\label{eq:lipschitz-condition}
		|f(x) - f(z)| \leq L\, |x - z| \qquad \forall\, x, z \in X.
	\end{equation}
	When $X = \R$, $f$ is \textit{globally} $L$-Lipschitz.
	When $X = [x - \varepsilon,\, x + \varepsilon]$ for some $\varepsilon > 0$, $f$ is \textit{locally} $L$-Lipschitz at $x$.
	The constant $L$ depends on $X$: locally Lipschitz everywhere does not imply globally Lipschitz ($f(x) = x^2$ is the standard counterexample).
\end{defn}

\begin{defn}[Continuity]\label{def:continuity}
	A function $f$ is \textit{continuous} at $x$ if
	\begin{equation}\label{eq:continuity}
		\forall\, \varepsilon > 0,\; \exists\, \delta > 0 :
		z \in [x - \delta,\, x + \delta]
		\;\Rightarrow\; |f(z) - f(x)| < \varepsilon.
	\end{equation}
	If $f$ is continuous at every $x \in \R$, we write $f \in C^0$.
	Many simple functions are continuous, and continuity is preserved under sums, products, composition, $\max$, and $\min$.
\end{defn}

\begin{thm}[Lipschitz implies continuity]\label{thm:lip-cont}
	If $f$ is locally Lipschitz at $x$, then $f$ is continuous at $x$.
\end{thm}
\begin{proof}
	Given $\varepsilon > 0$, choose $\delta = \varepsilon / L$.
	Then $|z - x| < \delta$ implies $|f(z) - f(x)| \leq L |z - x| < L \delta = \varepsilon$.
\end{proof}

The converse fails: $f(x) = \sqrt{|x|}$ is continuous everywhere but not Lipschitz on any interval containing the origin (it becomes ``infinitely steep'' as $x \to 0$).

\paragraph{Complexity of Lipschitz optimization}

Lipschitz continuity is the weakest useful regularity condition.
With it, optimization becomes possible --- but expensive.

\begin{thm}[Complexity for Lipschitz functions]\label{thm:lip-complexity}
	Let $f$ be $L$-Lipschitz on a bounded interval $[x_-, x_+]$ with diameter $D = x_+ - x_-$.
	A uniform sampling strategy --- evaluating $f$ at points spaced at most $2\varepsilon/L$ apart --- finds an $\varepsilon$-optimal solution in
	\begin{equation}\label{eq:lipschitz-complexity}
		\mathcal{O}\!\left(\frac{L\cdot D}{\varepsilon}\right)
	\end{equation}
	function evaluations.
	No algorithm can guarantee a better worst-case rate for the class of all $L$-Lipschitz functions \cite{hansen1995}.
\end{thm}

The bound scales as $1/\varepsilon$: each extra digit of accuracy costs ten times more evaluations.
Beyond $\varepsilon \approx 10^{-3}$ or $10^{-4}$, this becomes impractical.
And in higher dimensions, the situation is dramatically worse (we will see why in \Cref{sec:smooth-methods}).

Two practical complications make things harder still.
First, $L$ is typically unknown and hard to estimate --- yet algorithms need it.
Second, a method that works well for one class of functions may fail on another: the \textit{no free lunch theorem} \cite{serafino2014} says that, averaged over all possible objectives, every algorithm performs equally badly.
These limitations motivate the gradient-based methods of the next sections, which exploit differentiability to achieve rates like $\mathcal{O}(\log(1/\varepsilon))$.


% ────────────────────────────────────────────────────────────────
\subsection{Local Optimality}\label{ssec:local-opt}

Even if we stumble onto $x^*$, how do we \emph{recognize} it?
Verifying global optimality requires knowing $f^*$, which is the very thing we are trying to compute.
A weaker but tractable condition: local optimality.

\begin{defn}[Local minimum]\label{def:local-min}
	A point $x^*$ is a \textit{local minimum} of $f$ if there exists $\varepsilon > 0$ such that $f(x^*) \leq f(x)$ for all $x \in [x^* - \varepsilon,\, x^* + \varepsilon]$.
	A \textit{strict} local minimum additionally requires $f(x^*) < f(z)$ for all $z$ in that neighborhood other than $x^*$ itself.
\end{defn}

\begin{defn}[Strictly unimodal function]\label{def:unimodal}
	A function is \textit{strictly unimodal} on an interval $[x_-, x_+]$ if it has a unique minimum $x_*$ and is strictly decreasing on $[x_-, x_*]$ and strictly increasing on $[x_*, x_+]$.
\end{defn}

Most functions are not unimodal globally.
But near any local minimum $x^*$, a typical $C^1$ function \emph{is} unimodal within some neighborhood --- the \textit{attraction basin} of $x^*$.
Every local optimum has its own basin, and once we land inside one, finding the minimum is comparatively easy.
The trouble is that all basins ``look the same'' from the inside: a local method cannot tell whether the local minimum it finds is the global one.
Finding the right basin is an entirely different --- and much harder --- problem, which we defer to \Cref{ssec:global-opt}.

For now, we focus on finding \emph{some} local minimum efficiently.
If $f$ is strictly unimodal on $[x_-, x_+]$, placing two test points $x'_-$ and $x'_+$ inside the interval tells us which half contains the minimum:
if $f(x'_-) \geq f(x'_+)$, then $x^* \in [x'_-, x_+]$;
if $f(x'_-) \leq f(x'_+)$, then $x^* \in [x_-, x'_+]$ \cite{bazaraa2006}.
By iterating, we shrink the interval to any desired width $\delta$.
The question is: where should we place $x'_-$ and $x'_+$?


% ────────────────────────────────────────────────────────────────
\subsection{Golden Ratio Search}\label{ssec:golden-ratio}

The strategy is to optimize worst-case behavior: shrink the interval as quickly as possible, regardless of which half gets discarded.
Each iteration drops either $[x_-, x'_-]$ or $[x'_+, x_+]$; since we don't know which, these two segments should be equal.
This means choosing a ratio $r \in (1/2, 1)$ and placing $x'_- = x_- + (1-r)D$ and $x'_+ = x_- + rD$.

Smaller $r$ gives faster shrinkage, but there is a catch: if $r$ is close to $1/2$, the two test points nearly coincide, and in the next iteration both must be recomputed from scratch.
To save one evaluation per step, we need the surviving test point to become one of the two new test points.
This imposes $r : 1 = (1 - r) : r$, i.e., $r^2 = 1 - r$, giving
\begin{equation}\label{eq:golden-ratio}
	r = \frac{\sqrt{5} - 1}{2} \approx 0.618,
\end{equation}
the reciprocal of the golden ratio $g = (1 + \sqrt{5})/2 \approx 1.618$.

\begin{algorithm}
	\caption{Golden Ratio Search (GRS)}
	\label{alg:GRS}
	\begin{algorithmic}[1]
		\Require Function $f$, interval $[x_-, x_+]$, precision $\delta$
		\Ensure Approximate minimizer of $f$ in $[x_-, x_+]$
		\Procedure{GRS}{$f, x_-, x_+, \delta$}
			\State $(x'_-, x'_+) \gets (x_- + (1-r)(x_+ - x_-),\; x_- + r(x_+ - x_-))$
			\While{$x_+ - x_- > \delta$}
				\If{$f(x'_-) > f(x'_+)$}
					\State $(x_-, x'_-, x'_+) \gets (x'_-, x'_+, x_- + r(x_+ - x_-))$
				\Else
					\State $(x_+, x'_+, x'_-) \gets (x'_+, x'_-, x_- + (1-r)(x_+ - x_-))$
				\EndIf
			\EndWhile
		\EndProcedure
	\end{algorithmic}
\end{algorithm}

After $k$ iterations the interval has width $D r^k$, so the method stops when $k \approx 4.78 \log(D / \delta)$ --- exponentially fast in $\delta$, a huge improvement over the $\mathcal{O}(1/\varepsilon)$ of Lipschitz sampling.
With naive bisection ($r = 0.5$) and two evaluations per step, the constant would be $\approx 6.64$ instead of $4.78$.
Golden ratio search is asymptotically optimal among methods using only function values \cite{bazaraa2006}; a slightly better finite-$k$ variant uses Fibonacci numbers.

Note that $\delta$ and $\varepsilon$ are different: $\delta$ controls the interval width, not the gap $A(x)$.
If $f$ is $L$-Lipschitz, then $A(x^k) \leq L\delta$, so reaching $A \leq \varepsilon$ requires $k \approx 4.78 \log(LD / \varepsilon)$.


% ────────────────────────────────────────────────────────────────
\subsection{Derivatives and First-Order Models}\label{ssec:derivatives}

Golden ratio search uses only function values.
Can we do better if we also know \emph{in which direction} $f$ is decreasing?
For a linear function $f(x) = bx$, the sign of $b$ immediately tells us: left if $b > 0$, right if $b < 0$.
For a general nonlinear $f$, the derivative plays the same role locally.

\begin{defn}[Derivative and first-order model]\label{def:derivative}
	The \textit{derivative} of $f$ at $x$ is $f'(x) = \lim_{t \to 0} [f(x + t) - f(x)] / t$, provided the limit exists and is finite.
	When it does, the \textit{first-order model} $L_x(z) = f(x) + f'(x)(z - x)$ is the best linear approximation of $f$ near $x$.
	The sign of $f'(x)$ tells the local direction of decrease: $f'(x) < 0$ means $f$ is locally decreasing, $f'(x) > 0$ means increasing.
\end{defn}

If the left limit $f'_-(x) = \lim_{t \to 0^-}$ and right limit $f'_+(x) = \lim_{t \to 0^+}$ both exist but differ, $f$ is not differentiable at $x$.
The standard example: $f(x) = |x| = \max\{x, -x\}$, where $f'_-(0) = -1$ and $f'_+(0) = +1$.
Nondifferentiable functions arise naturally in practice ($\ell_1$-regularization, hinge loss); we deal with them in \Cref{sec:nonsmooth}.

Differentiability implies continuity (but not vice versa).
Derivatives of standard functions are known: $[x^k]' = kx^{k-1}$, $[e^x]' = e^x$, $[\ln x]' = 1/x$, $[\sin x]' = \cos x$, $[\cos x]' = -\sin x$.
The standard rules preserve differentiability:
$[\alpha f + \beta g]' = \alpha f' + \beta g'$ (linearity),
$[f \cdot g]' = f'g + fg'$ (product),
$[f/g]' = (f'g - fg')/g^2$ (quotient),
$[f(g(\cdot))]' = f'(g(\cdot)) \cdot g'$ (chain rule).
Two important operations that do \emph{not} preserve differentiability: $\max\{f, g\}$ and $\min\{f, g\}$.
In practice, automatic differentiation \cite{nocedal2006} computes derivatives efficiently, so computing them is not our bottleneck.

\paragraph{Regularity classes}

We write $f \in C^k$ when $f$ has $k$ continuous derivatives.
The hierarchy matters for optimization:
$f \in C^1$ means $f'$ is continuous, which implies $f$ is Lipschitz on every bounded interval (by the extreme value theorem applied to $|f'|$).
$f \in C^2$ means $f''$ exists and is continuous, giving us curvature information.
The best case for univariate optimization is $f \in C^3$: both $f$ and $f'$ are Lipschitz, and we get quadratic convergence from Newton's method.

\paragraph{Local optimality and derivatives}

The connection is immediate.
If $f'(x) \neq 0$, then $f$ is either increasing or decreasing at $x$, so $x$ cannot be a local minimum.
Therefore $f'(x^*) = 0$ at every local minimum --- a \textit{necessary} condition.
But $f'(x) = 0$ also holds at maxima and saddle points (inflection points with zero slope).
To distinguish them, we need the second derivative: $f''(x^*) > 0$ confirms a local minimum, $f''(x^*) < 0$ a local maximum, and $f''(x^*) = 0$ is inconclusive.
Finding local minima therefore reduces to finding roots of $f'$ --- and checking the sign of $f''$.


% ────────────────────────────────────────────────────────────────
\subsection{Dichotomic Search}\label{ssec:dichotomic}

When derivative information is available, we can find a root of $f'$ by bisection on its sign.
The intermediate value theorem guarantees: if $f' \in C^0$ with $f'(x_-) < 0$ and $f'(x_+) > 0$, then $\exists\, x \in [x_-, x_+]$ with $f'(x) = 0$.

\begin{algorithm}
	\caption{Dichotomic Search (DS)}
	\begin{algorithmic}[1]
		\Procedure{DS}{$f, x_-, x_+, \varepsilon$}
			\Repeat
				\State $x \gets (x_- + x_+) / 2$
				\If{$|f'(x)| \leq \varepsilon$} \Return $x$ \EndIf
				\If{$f'(x) < 0$}
					\State $x_- \gets x$ \Comment{Minimum is to the right}
				\Else
					\State $x_+ \gets x$ \Comment{Minimum is to the left}
				\EndIf
			\Until{$x_+ - x_- \leq \delta$}
		\EndProcedure
	\end{algorithmic}
\end{algorithm}

Each step halves the interval, giving linear convergence with $r = 0.5$.
The complexity is $k \approx 3.32 \log(D / \delta)$ --- fewer iterations than golden ratio search ($4.78$), though each iteration now requires a derivative evaluation instead of a function evaluation.
If $f'$ is $L$-Lipschitz (i.e., $f$ is $L$-smooth), the stopping criterion $|f'(x)| \leq \varepsilon$ relates to the interval width via $k \approx 3.32 \log(LD / (2\varepsilon))$.

\paragraph{Finding the initial interval}

What if we don't start with $f'(x_-) < 0 < f'(x_+)$?
A simple strategy: \textit{exponential expansion}.
Starting from a point with $f'(x_+) < 0$, double the step until the sign flips.

\begin{algorithm}
	\caption{Exponential Interval Expansion}
	\begin{algorithmic}[1]
		\State $\Delta x \gets 1$
		\While{$f'(x_+) \leq -\varepsilon$}
			\State $x_+ \gets x_+ + \Delta x$
			\State $\Delta x \gets 2\, \Delta x$ \Comment{Double the step}
		\EndWhile
	\end{algorithmic}
\end{algorithm}

If $f$ is unbounded below, the expansion never terminates --- a useful diagnostic.
If $f$ is coercive ($f(x) \to \infty$ as $|x| \to \infty$), it always finds a valid bracket.


% ────────────────────────────────────────────────────────────────
\subsection{Quadratic Interpolation}\label{ssec:quad-interpolation}

Plain bisection picks the midpoint, ignoring the \emph{magnitude} of $f'$ at the endpoints.
A better guess: fit a quadratic model $ax^2 + bx + c$ to the derivative information at $x_-$ and $x_+$, and use its minimum as the next iterate.

Matching $f'(x_-)$ and $f'(x_+)$ gives
\begin{equation}\label{eq:quad-interp-coefficients}
	a = \frac{f'(x_+) - f'(x_-)}{2(x_+ - x_-)},
	\qquad
	b = \frac{x_+ f'(x_-) - x_- f'(x_+)}{x_+ - x_-}.
\end{equation}
Setting $2ax + b = 0$ yields the interpolated estimate (the ``secant formula''):
\begin{equation}\label{eq:secant-formula}
	x = \frac{x_- f'(x_+) - x_+ f'(x_-)}{f'(x_+) - f'(x_-)}.
\end{equation}
This always falls between $x_-$ and $x_+$, but it can land dangerously close to an endpoint when one derivative is much larger than the other --- a case where the quadratic model is ``very skewed.''

The general lesson: \emph{never blindly trust a model}.
A safeguard clamps the estimate away from the boundaries:
\begin{equation}\label{eq:safeguard-clamping}
	x \gets \max\bigl\{x_- + \sigma(x_+ - x_-),\; \min(x_+ - \sigma(x_+ - x_-),\, x)\bigr\},
	\qquad \sigma \in (0, 0.5].
\end{equation}
In the worst case (model completely wrong), this degrades to linear convergence with $r = 1 - \sigma$.
When the model is good, convergence is superlinear.

\begin{thm}[Superlinear convergence of quadratic interpolation]\label{thm:secant-conv}
	If $f \in C^3$ with $f'(x^*) = 0$ and $f''(x^*) \neq 0$, then safeguarded quadratic interpolation converges with order $p = (1 + \sqrt{5})/2 \approx 1.618$ --- the golden ratio again --- provided the initial guess is close enough to $x^*$ \cite{sun2006}.
\end{thm}

Cubic interpolation, which uses all four pieces of information ($f$ and $f'$ at both endpoints), achieves quadratic convergence ($p = 2$) under the same hypotheses \cite{sun2006,nocedal2006}.
It is more complex to implement but often worthwhile in practice.


% ────────────────────────────────────────────────────────────────
\subsection{Newton's Method}\label{ssec:newton-univariate}

Newton's method pushes the approximation idea to its limit: instead of modelling $f$ linearly, it uses a \emph{quadratic} model, exploiting the second derivative.
Two equivalent views lead to the same update.

\textbf{View 1:} build a first-order model of $f'$ at $x^k$ and solve for its root.
\begin{equation}\label{eq:newton-view1}
	L'_k(x) = f'(x^k) + f''(x^k)(x - x^k) = 0
	\qquad \Longrightarrow \qquad
	x^{k+1} = x^k - \frac{f'(x^k)}{f''(x^k)}.
\end{equation}

\textbf{View 2:} build the second-order Taylor model of $f$ at $x^k$ and minimize it.
\begin{equation}\label{eq:newton-taylor-model}
	Q_k(x) = f(x^k) + f'(x^k)(x - x^k) + \frac{1}{2} f''(x^k)(x - x^k)^2.
\end{equation}
Setting $Q'_k(x) = 0$ yields the same update.

\begin{algorithm}
	\caption{Newton's Method (NM)}
	\begin{algorithmic}[1]
		\Procedure{NM}{$f, x, \varepsilon$}
			\While{$|f'(x)| > \varepsilon$}
				\State $x \gets x - f'(x) / f''(x)$
			\EndWhile
		\EndProcedure
	\end{algorithmic}
\end{algorithm}

\begin{thm}[Quadratic convergence of Newton's method]\label{thm:newton-conv}
	If $f \in C^3$ near $x^*$ with $f'(x^*) = 0$ and $f''(x^*) \neq 0$, and $x^0$ is sufficiently close to $x^*$, then Newton's method converges quadratically.
	By Taylor's theorem, there exists $\xi$ between $x^k$ and $x^*$ such that
	\begin{equation}\label{eq:newton-error}
		x^{k+1} - x^* = \frac{f'''(\xi)}{2\, f''(x^k)}\, (x^k - x^*)^2.
	\end{equation}
	Quadratic convergence means the number of correct digits roughly doubles at each step --- spectacular speed, but only when the initial guess is good enough.
\end{thm}

\begin{proof}
	Write $x^{k+1} - x^* = (x^k - x^*) - f'(x^k)/f''(x^k)$.
	Substituting $f'(x^*) = 0$ and applying Taylor's formula to $f'$ around $x^k$:
	$f'(x^*) - f'(x^k) - f''(x^k)(x^* - x^k) = \frac{1}{2} f'''(\xi)(x^* - x^k)^2$
	for some $\xi$ between $x^k$ and $x^*$.
	Dividing by $f''(x^k)$ gives the result.
	Since $f \in C^3$ and $f''(x^*) \neq 0$, there exists a neighborhood where $|f''(x)| \geq k_2 > 0$ and $|f'''(x)| \leq k_1 < \infty$.
	If $k_1 |x^k - x^*| / (2 k_2) < 1$, then $|x^{k+1} - x^*| < |x^k - x^*|$, and convergence follows by induction.
\end{proof}

Newton's method converges fast only locally.
Without a good starting point, it can diverge, oscillate, or converge to a maximum.
Global convergence requires safeguards (line search, trust region) --- topics we develop in the multivariate setting (\Cref{sec:smooth-methods}).
And $f''(x^k) = 0$ is trouble: the method is undefined, signaling a flat second-order model.


% ────────────────────────────────────────────────────────────────
\subsection{Global Optimization}\label{ssec:global-opt}

Local methods find local minima efficiently.
But what if the problem has many local minima and we need the best one?
Unless the objective has special structure, this is a fundamentally harder question.

\paragraph{Convexity: when local equals global}

\begin{thm}[Convexity and global optimality]\label{thm:convex-global}
	If $f$ is convex --- meaning $f''(x) \geq 0$ everywhere, or more generally $f(\alpha x + (1-\alpha)z) \leq \alpha f(x) + (1-\alpha) f(z)$ for all $x, z$ and $\alpha \in [0,1]$ --- then every local minimum is automatically a global minimum.
	Convexity ensures that $f'$ is nondecreasing, ruling out multiple stationary points (except along a flat segment).
\end{thm}

Convexity is the single most important structural assumption in optimization.
When it holds, local methods are already global --- we get the best of both worlds.
Many models are purposely constructed to be convex (SVM, ridge regression) precisely so that optimization is tractable.
``If you have the choice, choose convex'' \cite{boyd2008}.
We study convexity thoroughly in \Cref{sec:optimality}.

\paragraph{Spatial branch-and-bound}

When convexity doesn't hold and the global optimum is genuinely needed, the standard approach is \textit{spatial branch-and-bound}.
The idea: construct a convex lower approximation $\underline{f}(x) \leq f(x)$ on a bounded domain $X = [x_-, x_+]$, and solve $\min_x \underline{f}(x)$ to get a lower bound on $f^*$.

When the gap between the lower bound $\underline{f}(\bar{x})$ and the best known function value is too large, the domain is split into subregions.
Each subregion gets its own convex relaxation.
Subregions whose lower bound exceeds the best known value are \textit{pruned} --- they cannot contain the global minimum.

The method's efficiency hinges on the quality of $\underline{f}$.
A tight relaxation means aggressive pruning and fast convergence; a loose one means exploring many subregions.
Worst-case complexity is exponential --- keep dicing the domain until the pieces are tiny --- but practical performance depends on the degree of nonconvexity and the quality of the relaxation.

For Mixed-Integer Linear Programs, the branch-and-bound framework is especially effective: everything is ``trivial'' once the integer variables are fixed or relaxed.
Mixed-Integer Nonlinear Convex Programs are harder but still tractable.
General nonlinear nonconvex problems require additional machinery --- function transformations, convexification techniques --- typically implemented in specialized solvers.
These methods are immensely more efficient than blind search, yet immensely less efficient than local optimization.


% ────────────────────────────────────────────────────────────────
\subsection{Summary: Regularity and Speed}\label{ssec:univariate-summary}

The univariate methods form a hierarchy, each trading a stronger assumption for a faster rate:

\begin{table}[ht]
	\centering
	\begin{tabular}{llll}
		\toprule
		Method & Information used & Assumption & Complexity \\
		\midrule
		Uniform sampling & $f(x)$ only   & $L$-Lipschitz & $\mathcal{O}(LD/\varepsilon)$ \\
		Golden ratio     & $f(x)$ only   & unimodal      & $\mathcal{O}(\log(D/\delta))$ \\
		Dichotomic       & $f(x), f'(x)$ & $f \in C^1$, unimodal & $\mathcal{O}(\log(D/\delta))$ (faster constant) \\
		Quad.\ interpolation & $f'(x_\pm)$ & $f \in C^3$   & superlinear ($p \approx 1.618$) \\
		Newton           & $f'(x), f''(x)$ & $f \in C^3$ & quadratic ($p = 2$) \\
		\bottomrule
	\end{tabular}
	\caption{Univariate optimization methods: assumptions vs.\ speed.}
	\label{tab:univariate-methods}
\end{table}

These ideas --- models, line search, trust regions, regularity-dependent rates --- will reappear in every multivariate method.
But first, we need the multivariate calculus: gradients, Hessians, and the optimality conditions that generalize $f'(x^*) = 0$ and $f''(x^*) \geq 0$ to $\R^n$.